\documentclass[11pt,a4paper]{article}
\usepackage[utf8]{inputenc}
\usepackage[T1]{fontenc}
\usepackage[spanish]{babel}
\usepackage{lmodern}
\usepackage[a4paper,margin=2.2cm,headheight=14pt]{geometry}
\usepackage{booktabs,longtable,array,tabularx,enumitem,xcolor,hyperref,graphicx,fancyhdr,titlesec,amsmath,microtype}
\definecolor{azul}{HTML}{1F4E79}
\definecolor{azulclaro}{HTML}{EAF2F8}
\definecolor{gris}{HTML}{5B6573}
\hypersetup{colorlinks=true,linkcolor=azul,urlcolor=azul,citecolor=azul,pdftitle={Estructura del proyecto de análisis de datos},pdfauthor={Codex}}
\pagestyle{fancy}\fancyhf{}\lhead{Documentación del proyecto AED}\rhead{Estructura}\cfoot{\thepage}
\titleformat{\section}{\Large\bfseries\color{azul}}{\thesection}{0.6em}{}[\titlerule]
\titleformat{\subsection}{\large\bfseries\color{azul}}{\thesubsection}{0.6em}{}
\newcommand{\code}[1]{\texttt{\small #1}}
\newcolumntype{Y}{>{\raggedright\arraybackslash}X}
\begin{document}
\begin{titlepage}
\centering
\vspace*{1.8cm}
{\color{azul}\rule{\textwidth}{1.5pt}}\\[0.8cm]
{\Huge\bfseries Documentación de la estructura\\[0.15cm] del proyecto de análisis de datos\par}
\vspace{0.7cm}
{\Large Análisis Exploratorio de Datos provincial en España\par}
\vspace{1.4cm}
\begin{tabular}{rl}
\textbf{Proyecto:} & Estructura Productiva, Presión Inmobiliaria y su Impacto en el Ciclo Vital y Migratorio en España\\
\textbf{Autores del proyecto:} & Luis Molina Villar, Rubén Sánchez Munuera y Pablo Gallart Carbonell\\
\textbf{Institución:} & Universitat de València\\
\textbf{Fecha de documentación:} & 8 de noviembre de 2026
\end{tabular}
\vfill
\begin{minipage}{0.85\textwidth}\small\centering
Este documento describe la arquitectura, los ficheros, el flujo de transformación y los controles de calidad del proyecto entregado. La documentación se ha construido a partir de los archivos locales del repositorio, del informe PDF y de los datos incluidos en la carpeta \code{data/}.
\end{minipage}
\vfill
{\color{azul}\rule{\textwidth}{1.5pt}}
\end{titlepage}

\tableofcontents
\newpage

\section{Resumen ejecutivo}
El proyecto realiza un Análisis Exploratorio de Datos (AED) sobre las 50 provincias españolas peninsulares e insulares, excluyendo Ceuta y Melilla. El objetivo es estudiar si la estructura productiva provincial se relaciona de forma homogénea o heterogénea con la renta, la vivienda, la demografía y la migración.

El flujo completo es:
\begin{center}
\colorbox{azulclaro}{\parbox{0.92\textwidth}{\centering\textbf{fuentes INE/MIVAU} $\rightarrow$ \textbf{lectura española} $\rightarrow$ \textbf{normalización} $\rightarrow$ \textbf{unión por provincia-año} $\rightarrow$ \textbf{imputación} $\rightarrow$ \textbf{análisis 2019} $\rightarrow$ \textbf{correlaciones y clustering}}}
\end{center}

La unidad final de observación es una provincia en un año. El panel socioeconómico común cubre 2015--2021 y el análisis principal selecciona 2019, año anterior a la pandemia de COVID-19 y con disponibilidad común de todas las fuentes.

\section{Inventario de archivos}
\subsection{Archivos principales}
\begin{longtable}{p{0.29\textwidth}p{0.18\textwidth}p{0.45\textwidth}}
\toprule
\textbf{Archivo} & \textbf{Tipo} & \textbf{Responsabilidad}\
\midrule
\endhead
\code{ProyectoAED.Rmd} & R Markdown & Informe reproducible: planteamiento, preparación, EDA, correlaciones, clustering y conclusiones.\\
\code{Utilidades.Rmd} & R Markdown & Funciones auxiliares, catálogos de provincias, lectura, validación, imputación y gráficos.\\
\code{codigo\_extraido.R} & R & Versión ejecutable del código de \code{Utilidades.Rmd}, creada con \code{knitr::purl()}.\\
\code{instalar\_paquetes.R} & R & Script auxiliar de instalación/configuración disponible en el proyecto.\\
\code{ProyectoAED.tex} & LaTeX & Fuente LaTeX del informe principal.\\
\code{Proyecto\_AED.pdf} & PDF & Informe final original, de 10 páginas.\\
\code{ProyectoAED.Rproj} & RStudio & Configuración del proyecto RStudio.\\
\code{README.txt} & Texto & Nota de uso y contexto del repositorio.\\
\code{EnlaceRepo.txt} & Texto & Referencia al repositorio del proyecto.\\
\bottomrule
\end{longtable}

\subsection{Estructura de carpetas}
\begin{verbatim}
Entrega_ProyectoAED/
|-- ProyectoAED.Rmd              # informe reproducible
|-- Utilidades.Rmd               # funciones auxiliares
|-- codigo_extraido.R            # utilidades extraidas
|-- Proyecto_AED.pdf             # informe renderizado
|-- ProyectoAED.tex              # fuente LaTeX del informe
|-- data/
|   |-- raw/                     # ocho fuentes CSV originales
|   |-- Datos_SocioEcon.csv      # panel socioeconómico procesado
|   |-- Datos_SocioEcon.Rdata    # mismo panel en formato R
|   |-- Datos_Laborales.csv      # panel laboral procesado
|   `-- Datos_Laborales.Rdata    # mismo panel en formato R
|-- ProyectoAED_files/figure-latex/ # figuras del informe
|-- Definitions/                 # clase MDPI, bibliografía y logos
`-- output/pdf/                  # documentación generada
\end{verbatim}

\section{Fuentes de datos y dimensiones}
Se emplean ocho ficheros CSV de INE y MIVAU. Los ficheros están delimitados de forma distinta; por ello la lectura no se puede resolver con una única llamada genérica sin parámetros.

\begin{longtable}{p{0.26\textwidth}p{0.17\textwidth}p{0.15\textwidth}p{0.32\textwidth}}
\toprule
\textbf{Fuente local} & \textbf{Origen} & \textbf{Filas aprox.} & \textbf{Contenido y tratamiento}\
\midrule
\endhead
\code{coste\_vivienda.csv} & MIVAU & 13.711 & Precio por m$^2$, trimestre y régimen; se conserva vivienda libre y se promedia por provincia-año.\\
\code{distribucion\_poblacional.csv} & INE & 4.134 & Habitantes por provincia, periodo y sexo; se conserva \code{Sexo == "Total"}.\\
\code{edad\_maternidad.csv} & INE & 13.250 & Edad media; se conserva \code{Orden de nacimiento == "Todos"}.\\
\code{migracion\_interior.csv} & INE & 200.928 & Saldo por edad y sexo; se conserva ambos sexos y edades 19--35, sumando por provincia-año.\\
\code{Ocupados\_Sector\_Provincia.csv} & INE & 19.875 & Ocupación por ramas; construcción e industria se agrupan en secundario.\\
\code{pib\_per\_capita.csv} & INE & 4.914 & PIB por magnitud y territorio; se eliminan totales y se resuelven CCAA uniprovinciales.\\
\code{renta\_bruta\_media.csv} & INE & 4.920 & Renta bruta media por persona; se filtra el indicador exacto y se normalizan nombres.\\
\code{valor\_bruto\_sector.csv} & INE & 12.480 & VAB por rama; se agregan ramas en primario, secundario y terciario.\\
\bottomrule
\end{longtable}

\subsection{Dimensiones conceptuales}
\begin{itemize}[leftmargin=1.5em]
\item \textbf{Renta y riqueza:} \code{PIB\_per\_capita}, \code{Renta\_bruta\_media}.
\item \textbf{Mercado inmobiliario:} \code{Precio\_m2}.
\item \textbf{Demografía y ciclo vital:} \code{Habitantes}, \code{Edad\_Media\_Maternidad}, \code{Migracion\_Interior}.
\item \textbf{Estructura productiva:} ocupación y VAB de los sectores primario, secundario y terciario.
\end{itemize}

\section{Arquitectura del flujo de datos}
\subsection{Etapa 1: configuración y catálogo territorial}
El bloque de configuración carga los paquetes requeridos y establece opciones globales de \code{knitr}: código visible, advertencias y mensajes ocultos, figuras centradas, tamaño de 6 por 4 pulgadas y ancho máximo del texto del PDF. \code{Utilidades.Rmd} construye además \code{ine\_provincias}, una tabla de 50 códigos INE y nombres oficiales, y \code{ine\_provincias\_join}, que añade una clave normalizada para unir fuentes con tildes, barras y nombres alternativos.

\subsection{Etapa 2: descarga y lectura}
\code{comprobacion\_descarga\_datos()} crea \code{data/raw} si es necesario, comprueba la existencia de los ocho archivos y descarga únicamente los que no existen o tienen tamaño cero. \code{leer\_archivo\_ine()} aplica el delimitador, la configuración decimal/millar y las etiquetas de ausentes \code{""}, \code{"NA"}, \code{"null"} y \code{".."}. \code{leer\_datos()} carga cada fuente en un objeto con sufijo \code{\_raw}.

\subsection{Etapa 3: normalización y claves}
Cada tabla se lleva a formato tidy. Se separa el texto \code{"código provincia nombre"}, se rellenan códigos a dos dígitos y se genera:
\[
\code{id} = \code{año\_provincia},\qquad \text{por ejemplo, }\code{2019\_28}.
\]
La clave es común a todas las tablas socioeconómicas; para los datos laborales se utiliza \code{(id, Sector)}.

\subsection{Etapa 4: agregación de fuentes}
\begin{enumerate}[leftmargin=1.7em]
\item Vivienda: se corrigen nombres de columnas, se recuperan las comunidades uniprovinciales, se filtra \code{Regimen == "Libre"} y se calcula la media trimestral.
\item Población: se conserva el total de ambos sexos y se renombra \code{Total} como \code{Habitantes}.
\item Maternidad: se conserva el total de órdenes de nacimiento y se convierte el indicador numérico con configuración española.
\item Migración: se conserva \code{Ambos sexos}, edades 19--35 y se suman los saldos de esas edades.
\item Ocupación: se pasa a columnas, se calcula \code{Secundario = Industria + Construccion} y se vuelve a formato largo.
\item PIB y renta: se filtran indicadores, se corrigen nombres y se resuelven duplicidades entre provincia y comunidad autónoma uniprovincial.
\item VAB: se excluyen impuestos y PIB a precios de mercado; las ramas se suman por sector.
\end{enumerate}

\subsection{Etapa 5: paneles y persistencia}
El panel socioeconómico se forma con seis \code{inner\_join()} por \code{id}. El panel laboral se forma con un cruce por \code{id} y \code{Sector}. Después los sectores se pasan a columnas, se unen ambos bloques y se separa la clave en \code{Periodo} e \code{id\_prov}. El panel final se guarda en CSV delimitado por punto y coma y en formato binario RData.

\section{Análisis exploratorio}
\subsection{Datos perdidos}
Se restringe el panel a 2015--2021. Se contabilizan ausentes por variable y se listan las filas con ausentes. La imputación solo se aplica a \code{Renta\_bruta\_media} y \code{Personas\_Ocupadas\_Primario}, por provincia y en orden temporal: media entre el primer valor válido anterior y posterior; en los extremos se usa el único vecino disponible.

\subsection{Análisis univariante}
Para 2019 se calculan participaciones nacionales de ocupación y VAB. Se muestran histogramas con media, mediana, desviación típica, curtosis y atípicos a más de tres desviaciones. Los boxplots usan la regla $1.5\times IQR$, con transformaciones visuales para población y migración, pero calculando atípicos sobre la escala original.

\subsection{Análisis bivariante y multivariante}
Se derivan cuatro indicadores: densidad, esfuerzo de vivienda, migración joven por 1.000 habitantes y brecha renta/PIB. La matriz global usa Spearman y ordenamiento jerárquico por $1-|\rho|$. Para explorar heterogeneidad, se estandarizan 11 variables, se calcula distancia euclídea, se aplica \code{hclust(method = "ward.D2")} y se cortan cuatro grupos. Finalmente se cartografían las provincias con \code{mapSpain} y se calculan matrices de correlación por cluster.

\section{Controles de calidad}
El proyecto incluye controles en puntos críticos:
\begin{itemize}[leftmargin=1.5em]
\item \code{validar\_clave\_unica()} detiene la ejecución ante duplicados de la clave esperada.
\item \code{validar\_provincias\_50()} detecta provincias fuera del alcance y avisa si falta alguna de las 50.
\item Se usan \code{inner\_join()} para que el panel final solo contenga instancias comunes a todas las fuentes.
\item Se eliminan explícitamente Ceuta y Melilla del análisis cartográfico y del clustering.
\item Se conservan dos formatos de salida: CSV auditable y RData reproducible.
\end{itemize}

\section{Productos generados por el proyecto}
\begin{longtable}{p{0.36\textwidth}p{0.5\textwidth}}
\toprule\textbf{Producto} & \textbf{Uso}\
\midrule
\endhead
\code{data/Datos\_SocioEcon.csv/.Rdata} & Variables de vivienda, población, maternidad, migración, PIB y renta por provincia-año.\\
\code{data/Datos\_Laborales.csv/.Rdata} & Personas ocupadas y VAB por provincia-año-sector.\\
\code{Proyecto\_AED.pdf} & Informe científico con metodología, figuras, resultados y conclusiones.\\
\code{ProyectoAED\_files/figure-latex/} & Figuras incrustadas o auxiliares del informe.\\
\code{output/pdf/} & Esta documentación: estructura, referencia de R y guion oral.\\
\bottomrule
\end{longtable}

\section{Reproducción recomendada}
\begin{enumerate}[leftmargin=1.7em]
\item Abrir el proyecto desde \code{ProyectoAED.Rproj} para conservar las rutas relativas.
\item Verificar que R y los paquetes listados en la documentación de R están disponibles.
\item Ejecutar primero \code{Utilidades.Rmd} o extraerlo con \code{knitr::purl()}.
\item Ejecutar \code{ProyectoAED.Rmd} desde la raíz del proyecto, manteniendo \code{data/raw}.
\item Revisar que las validaciones no devuelvan duplicados y que el periodo común siga siendo 2015--2021.
\item Renderizar el informe y comprobar las figuras y tablas antes de entregar.
\end{enumerate}

\section{Limitaciones documentadas}
El análisis es exploratorio y provincial: no prueba causalidad. Las participaciones laborales y de VAB están expresadas respecto al total nacional, por lo que el tamaño de la provincia influye. El clustering depende de las variables elegidas, la estandarización, la distancia y el método Ward.D2. Los clusters 3 y 4 tienen muy pocas provincias, de modo que sus correlaciones internas deben interpretarse con especial cautela. El año 2019 es una fotografía pre-pandemia; extender el análisis a 2015--2021 permitiría estudiar la evolución temporal.

\section*{Cierre}
La estructura del proyecto separa claramente datos brutos, funciones, análisis y productos. Esa separación permite auditar el origen de cada indicador, repetir el pipeline y distinguir entre decisiones de preparación de datos y conclusiones substantivas.

\end{document}
